\chapter{总结与展望}
\label{chap:conclusion}

\section{全文总结}
\label{sec:conclusion-summary}

本文围绕建筑幕墙韧性评估中的现实痛点，设计并实现了一套基于 AI Agent 的软件系统。本文工作的重点不是重新提出某个检测模型，而是解决现实工程中更关键的链路问题：如何把海量幕墙图像、不同原子检测能力、人工复核信息和项目级报告生成组织为一个完整、可追踪、可部署的系统，并通过 Agent 提升任务调度和信息整合能力。

全文首先分析了建筑幕墙巡检从人工经验向智能评估转型过程中面临的痛点，指出现有单点检测方法难以解决项目级图像管理和报告生成问题。随后，本文梳理了建筑立面检测、计算机视觉和 AI Agent 相关研究，说明 Agent 工具调用和多阶段工作流为幕墙评估提供了新的技术机会。在问题建模基础上，本文提出受控 Agent 调度和分层信息整合策略，将分类 Agent、图像级总结 Agent 和项目级报告 Agent 分别嵌入检测前、检测后和项目报告阶段。最后，本文完成了包含前端、API、BIZ、Activity 服务、数据库、对象存储、消息队列和部署配置的系统原型。

\section{主要成果}
\label{sec:conclusion-results}

本文取得的主要成果包括以下三点。第一，设计了面向多材质幕墙图像的 Agent 路由机制，使系统能够根据图像内容动态决定需要执行的检测能力。该机制避免了固定执行全部模型的低效方式，也使检测流程更符合幕墙材料差异。

第二，提出并实现了分层信息整合策略。系统先保存原子检测结构化结果，再生成图像级总结，最后生成项目级报告。该策略解决了多图像、多任务、多区域信息直接汇总导致的上下文压力和信息噪声问题，使报告生成过程更稳定、更可解释。

第三，构建了完整工程闭环。系统通过服务化架构、gRPC 协议、异步 Worker、回调接口、数据库状态机、对象存储、消息队列和 WebSocket 推送，使 Agent 能力从一次性调用变成可运行、可追踪、可恢复的业务流程。这一点体现了本文作为软件工程毕业设计的工作量和系统复杂度。

从毕业设计工作量角度看，本文完成的不只是若干独立功能点，而是一条跨前端、后端、模型服务、对象存储、消息队列和 Agent 平台的完整链路。系统中任何一个环节缺失，最终用户都无法获得稳定的智能评估体验。因此，本文的主要成果既包括方法设计，也包括对方法的工程落地验证。

\section{不足之处}
\label{sec:conclusion-limitations}

本文仍存在不足。首先，五类原子检测能力虽然已经接入系统，但模型本身不是本文重点，尚未围绕真实大规模幕墙数据集进行系统精度评估。其次，Agent 的领域知识主要来自提示词和结构化输入，尚未接入可检索的幕墙规范、专家案例和知识库。再次，报告生成质量仍需要通过更多真实项目和专家评价进行校准。最后，系统作为毕业设计原型，在监控告警、权限精细化、分布式扩展、模型版本管理和报告导出方面仍有继续完善空间。

\section{未来工作}
\label{sec:conclusion-future}

未来工作可以从四个方向展开。第一，扩充建筑幕墙真实场景数据集，对原子检测能力进行更严格的评估和优化。第二，引入检索增强生成机制，将幕墙工程规范、维护标准和专家案例作为 Agent 的外部知识来源，提高报告专业性。第三，完善系统工程能力，包括分布式 Worker、任务优先级、模型版本切换、监控告警和数据备份。第四，探索与 BIM 平台、无人机巡检平台和建筑运维系统对接，使图像检测结果能够进一步映射到建筑空间模型中。

总体而言，本文验证了 AI Agent 在建筑幕墙韧性评估中的应用价值。Agent 不是替代工程系统，而是通过受控调度和分层信息整合补足传统工程系统在语义判断和报告生成方面的不足。本文系统为建筑幕墙智能运维提供了一条可落地的技术路线，也为 AI Agent 在垂直工程软件中的应用提供了参考。
